\lecture{9}{24 March.}{}
\begin{eg}[Coupon collector problem]
    Every purchase we make, we get a coupon, which is uniformly randomly chosen from one of \(n\) types, independent of all previous coupons. The goal is to complete a set: have at least one of every type of coupon. We can define a random variable \(T\), which is the number of purchases before we have a full set. Then, what is the distribution of \(T\)? i.e. for \(k \in \mathbb{N} \), what is \(p(k) = \mathbb{P} (T = k)\)?    
\end{eg}

\begin{explanation}
    Suppose 
    \[
        S = \left\{ (s_1, s_2, \dots) : s_i \in [n] \text{ is the type of the } i\text{-th coupon}  \right\}. 
    \]
    Then \(\mathbb{P} (s_j = j) = \frac{1}{n}\) since coupons are uniform. By independence, 
    \[
        \mathbb{P} (\left\{ s_i = j_i : 1 \le i \le k \right\} ) = \frac{1}{n^k} \text{ for any } j_i \in [n]. 
    \] 
    Also,
    \[
        T((s_1, s_2, \dots )) = \begin{dcases}
            \min \left\{ k: (s_1, \dots , s_k) \text{ contain all type form}  \right\} , &\text{ if such } k \text{ exists};\\
            -1, &\text{ if such } k \text{ does not exist}.
        \end{dcases}
    \]
    Thus, \(\mathbb{P} (T = k)\) for \(k < n\) is \(0\) since we need at least \(n\) purchases to get \(n\) coupons. Now since the first \(n\) coupons are uniformly distributed over all possible sequences. Thus, 
    \begin{align*}
        \# \text{ of possible sequences } &= n^n \text{ (product rule)} \\
        \# \text{ of good sequences } &= n! \text{ (permutation of all types of coupon)}.  
    \end{align*}
    Thus, 
    \[
        \mathbb{P} (T = n) = \frac{n!}{n^n} \thickapprox \left( \frac{1}{e} \right)^n. 
    \] 
    Now we want to compute \(\mathbb{P} (T = k)\) for \(k > n\). We can first consider \(\mathbb{P} (T > k)\). Let 
    \[
        E_i^{(k)} = \left\{ \text{coupon } i \text{ missing from the first } k \text{ coupons}    \right\}. 
    \]  
    Then,
    \[
        \left\{ T > k \right\} \coloneqq \bigcup_{i=1}^{n} E_i^{(k)}.  
    \]
    Note that for all \(i\) 
    \[
        \mathbb{P} \left( E_i^{(k)} \right) = \left( \frac{n-1}{n} \right)^k.  
    \]
    \begin{question}
        Are \(E_1^{(k)}, E_2^{(k)}, \dots , E_n^{(k)}\) independent? 
    \end{question}
    \begin{answer}
        Since \(\mathbb{P} \left( E_i^{(k)} \right) = \left( \frac{n-1}{n} \right)^k\) for all \(i\), so
        \[
            \prod _{i=1}^n \mathbb{P} (E_i^{(k)}) = \left( \frac{n-1}{n} \right)^{kn},
        \] but
        \[
            \mathbb{P} \left( \bigcap_{i=1}^{n} \mathbb{P} (E_i^{(k)})  \right) = 0, 
        \]
        so the answer is no.
    \end{answer}
    By inclusion-exclusion principle,
    \begin{align*}
        \mathbb{P} \left( \bigcup_{i=1}^{n} E_i^{(k)}  \right) &= \sum_{\varnothing \neq I \subseteq [n]} (-1)^{\vert I \vert + 1} \mathbb{P} \left( \bigcap_{i \in I} E_i^{(k)}  \right) = \sum_{\varnothing \neq I \subseteq [n]} (-1)^{\vert I \vert  + 1} \left( \frac{n - \vert I \vert }{n} \right)^k \\
        &= \sum_{r=1}^{n-1} \binom{n}{r} (-1)^{r+1} \left( \frac{n-r}{n} \right)^k. 
    \end{align*}
    Thus, 
    \begin{align*}
        \mathbb{P} (T = k) &= \mathbb{P} (T > k - 1) - \mathbb{P} (T > k) \\
        &= \sum_{r=1}^{n-1} \binom{n}{r} (-1)^{r+1} \left( \frac{n-r}{n} \right)^{k-1} - \sum_{r=1}^{n-1} \binom{n}{r} (-1)^{r+1} \left( \frac{n-r}{n} \right)^k \\
        &= \sum_{r=1}^{n-1} \binom{n}{r} (-1)^{r+1} \left( \frac{n-r}{n} \right)^{k-1} \frac{r}{n}.    
    \end{align*}
    

\end{explanation}

\begin{definition}[Expected value]
    Let \((S, \mathbb{P})\) be a probability space and let \(X: S \to \mathbb{R} \) be a discrete random variable. Then the expected value of \(X\) 
    \[
        \mathbb{E} [X] = \sum_{x_i} x_i p(x_i) = \sum_{x_i} x_i \mathbb{P} (X = x_i).  
    \]  
\end{definition}

\begin{remark}
    Expected value is what we expect to see on average if we repeat the random experiment many times.
\end{remark}

Suppose we repeat \(N\) times, and let \(n_i\) to be the number of times we see \(X = x_i\), then the average value of \(X\) is 
\[
    \frac{\sum_{i} x_i n_i }{N} = \sum_{i} x_i \frac{n_i}{N} \thickapprox \sum_{i} x_i \mathbb{P} (X = x_i) \text{ when } N \to \infty .   
\]    

\begin{eg}
    Two multiple choice questions, each with four options. If we choose answers uniformly at random, how many do we expect to get correct?
\end{eg}
\begin{explanation}
    Let \(X = \#\) of correct answers. Then, 
    \renewcommand{\arraystretch}{1.5}
    \begin{table}[H]
        \centering
        \begin{tabular}{c|c}
            \toprule
                \(x_i\)  & \(\mathbb{P} (X = x_i)\)   \\
            \midrule
                \(0\)  & \(\left( \frac{3}{4} \right)^2 \)  \\

                \(1\)  & \(2 \left( \frac{1}{4} \right) \left( \frac{3}{4} \right)  \)   \\
                \(2\)  & \(\left( \frac{1}{4} \right)^2 \)   \\
            \bottomrule
        \end{tabular}
    \end{table} 
    Thus, 
    \[
        \mathbb{E} [X] = \sum_{i} x_i \mathbb{P} (X = x_i) = 0 \cdot \frac{9}{16} + 1 \cdot \frac{6}{16} + 2 \cdot \frac{1}{16} = \frac{1}{3}. 
    \]
\end{explanation}

\begin{claim}
    \[
        \mathbb{E} [X] = \sum_{s \in S} X(s) \mathbb{P} (\left\{ s \right\} ) 
    \]
\end{claim}
\begin{explanation}
        \begin{align*}
            \mathbb{E} [X] &= \sum_{x_i} x_i \mathbb{P} (X = x_i) = \sum_{x_i} x_i \mathbb{P} (\left\{ s : X(s) = x_i \right\} ) \\
            &= \sum_{x_i} x_i \mathbb{P} \left( \bigcup_{s: X(s) = x_i} \left\{ s \right\}   \right) = \sum_{x_i} x_i \sum_{\substack{s \in S \\ X(s) = x_i}} \mathbb{P} (\left\{ s \right\} ) = \sum_{s \in S} X(s) \mathbb{P} \left( \left\{ s \right\}  \right).    
        \end{align*}
\end{explanation}

\section{Expected value of functions of random variable}
Let \((S, \mathbb{P} )\) be a probability space. Let \(X: S \to \mathbb{R} \) be a random variable. Let \(g: \mathbb{R} \to \mathbb{R} \) be a function. Then, \(g(X) = g \circ X: S \to \mathbb{R} \) is another random variable.

\begin{claim}
    \[
        \mathbb{E} [g(X)] = \sum_{x_i} g(x_i) \mathbb{P} (X = x_i). 
    \]
\end{claim}

\begin{explanation}
    By the earlier claim, 
    \begin{align*}
        E [g(X)] &= \sum_{s \in S} g(X(s)) \mathbb{P} \left( \left\{ s \right\}  \right) = \sum_{x_i} \sum_{\substack{s \in S \\ X(s) = x_i}} g(X(s)) \mathbb{P} \left( \left\{ s \right\}  \right) \\
        &= \sum_{x_i} g(x_i) \sum_{\substack{s \in S \\ X(s) = x_i}} \mathbb{P} \left( \left\{ s \right\}  \right) = \sum_{x_i} g(x_i) \mathbb{P} (X = x_i).       
    \end{align*}
\end{explanation}

\begin{proposition}
    Let \(X_1, X_2, \dots , X_n: S \to \mathbb{R} \) be random variables on a probability space and \(\alpha _1, \alpha _2, \dots , \alpha _n \in \mathbb{R} \). Then if 
    \[
        X = \sum_{i=1}^n \alpha _i X_i : S \to \mathbb{R} , \text{ we have } \mathbb{E} [X] = \sum_{i=1}^n \alpha _i \mathbb{E} [X_i].   
    \] 
    This is called linearity of expectation. 
\end{proposition}

\begin{proof}
    By earlier claim,
    \begin{align*}
        \mathbb{E} [X] &= \sum_{s \in S} X(s) \mathbb{P} \left( \left\{ s \right\}  \right) = \sum_{s \in S} \left( \sum_{i=1}^n \alpha _i X_i(s)  \right) \mathbb{P} \left( \left\{ s \right\}  \right) \\
        &= \sum_{s \in S} \sum_{i=1}^n \alpha _i X_i(s) \mathbb{P} \left( \left\{ s \right\}  \right) = \sum_{i=1}^n \alpha _i \sum_{s \in S} X_i(s) \mathbb{P} \left( \left\{ s \right\}  \right) \\
        &= \sum_{i=1}^n \alpha _i \mathbb{E} [X_i].            
    \end{align*}
\end{proof}

\begin{remark}
    We do not require any independence.
\end{remark}

\begin{definition}
    Random variables \(X_1, X_2, \dots \) are independent if for any \(x_1, x_2, \dots \), the events 
    \[
        E_i = \left\{ X_i = x_i \right\} \text{ for all } i 
    \] 
    are independent.
\end{definition}

\begin{eg}
    Multiple choice: each question has \(k\) options, one of each which is correct. Suppose there are \(n\) questions in total. If we guess uniformly at random, independently, what is the expected score?  
\end{eg}

\begin{answer}      
    Let \(X\) be the number of correct numbers. Then \(X\) takes values in \([n] \cup \left\{ 0 \right\} \). For each \(0 \le r \le n\), 
    \[
        \mathbb{P} (X = r) = \binom{n}{r} \left( \frac{1}{k} \right)^r \left( \frac{k-1}{k} \right)^{n-r}.  
    \]
    Thus, 
    \[
        \mathbb{E} [X] = \sum_{x_i} x_i \mathbb{P} (X = x_i) = \sum_{r = 0}^n r \mathbb{P} (X = r) = \sum_{r = 1}^n r \binom{n}{r} \left( \frac{1}{k} \right)^r \left( \frac{k-1}{k} \right)^{n-r} = \frac{n}{k}.   
    \]    
\end{answer}        

\begin{answer}
    Let 
    \[
        X_i = \begin{dcases}
            1, &\text{ if } i \text{-th question is correct}  ;\\
            0, &\text{ if } i \text{-th question is not correct} ,
        \end{dcases}
    \]
    which is the indicator random variable for the event of getting \(i\)-th question correct. Then,
    \begin{align*}
        \mathbb{E} [X_i] &= 0 \cdot \mathbb{P} (i\text{-th question is not correct} ) + 1 \cdot \mathbb{P} (i\text{-th question is correct} ) \\
        &= 0 \cdot \frac{k-1}{k} + 1 \cdot \frac{1}{k} = \frac{1}{k}.
    \end{align*}
    Then, the number of correct answers is given by 
    \[
        X = \sum_{i=1}^n X_i. 
    \]
    By linearity of expectation,
    \[
        \mathbb{E} [X] = \sum_{i=1}^n \mathbb{E} [X_i] = \frac{n}{k}. 
    \]
\end{answer}

\begin{question}
    Three investment opportunities: 
    \begin{itemize}
        \item [1.] Guaranteed to get \(\$ 1\) after a year. 
        \item [2.] We get \(\$ 1000\) with probability \(\frac{1}{1000}\) and \(\$ 0\) at probability \(0\). 
        \item [3.] Get \(\$ \frac{2000}{999}\) with probability \(\frac{999}{1000}\) and \(\$ -1000\) at probability \(\frac{1}{1000}\).        
    \end{itemize}
    Which should we choose?
\end{question}